% Generated by scripts/contract.py; longtable body only.
\cref{lem:tc-ground-pins} & \lean{DeciENSSE.Paper.lem_tc_ground_pins} & The recursive list groundPins enumerates constructor and extreme-leaf positions; a Boolean selects lower or upper bounds, and PathPin.Holds interprets their observations. The equivalence also holds at arity zero. \\
\cref{lem:tc-ranked} & \lean{DeciENSSE.Paper.lem_tc_ranked} & Finite variable sets are initial segments; Ranked.Extension packages the embedding, source/sink partition, and rank assumptions. Old ranks need not be zero; the result also holds at arity zero. \\
\cref{lem:mo-transport} & \lean{DeciENSSE.Paper.lem_mo_transport} & The expansion is expand l v' with block cuts cutP l v' i; only the data of the alignment lemma are assumed, from which (x, y) is a comparison; letters are optional list entries. Marker avoidance is not required. \\
\cref{lem:mo-types} & \lean{DeciENSSE.Paper.lem_mo_types} & The number of types is the cardinality of the type space, the codomain of anchoredType at the protected letters; the bound holds at every level below the depth of every nonempty word. Level zero is included; Inhabited supplies the default letter of the canonical hierarchy, which is unused because the word and its hierarchy levels are nonempty. \\
\cref{lem:mo-exchange} & \lean{DeciENSSE.Paper.lem_mo_exchange} & The letter map h is identify a b; dictionaries are compared through alphCard and start counts through sCount. The canonical hierarchy uses an Inhabited alphabet; a letter is available from the nonempty hole, and the default is never used. \\
\cref{lem:mo-dictionary}(1) & \lean{DeciENSSE.Paper.lem_mo_dictionary_1} & Anchored normality is PosNormal; dictionaries are compared through alphCard and start counts through sCount. The canonical hierarchy uses an Inhabited alphabet; a letter is available from the nonempty hole, and the default is never used. \\
\cref{lem:mo-dictionary}(2) & \lean{DeciENSSE.Paper.lem_mo_dictionary_2} & g\_i is alphCard w i and b\_i is the cardinality of boundaryPairs; A-sharp and G-sharp are anchoredAllowance and anchoredDictionaryBound. The unnumbered consequence is included. Inhabited supplies the default letter of the canonical hierarchy, which is unused on nonempty holes. \\
\cref{lem:mo-unary} & \lean{DeciENSSE.Paper.lem_mo_unary} & The unary word z\textasciicircum{}m is List.replicate m z. \\
\cref{lem:mo-agreement} & \lean{DeciENSSE.Paper.lem_mo_agreement} & Agreement with channels is PosAgree, between monitors with a shared core type and possibly different alphabets; the support and the family of blocks are lists. \\
\cref{lem:mo-summary}(1) & \lean{DeciENSSE.Paper.lem_mo_summary_1} & A column word is a list of columns, maps from Fin k to Option Gamma with none as padding, and its current tuple is dec; the fields form a PosColumnState and their meanings are psumm; the prefixes of a valid word are its takes. State finiteness and marker avoidance are not required for this invariant. \\
\cref{lem:mo-summary}(2) & \lean{DeciENSSE.Paper.lem_mo_summary_2} & The final positive summary is psumm of the tuple, by item (1); determination is stated as equal summaries giving equal rows of the derived channel, followed by the explicit formulas; the first-letter field is an optional Prop. State finiteness and marker avoidance are not required for these formulas. \\
\cref{lem:mo-columns} & \lean{DeciENSSE.Paper.lem_mo_columns} & The family is a list of at most k blocks; subsequences are sublists, and injectivity is stated on the list. \\
\cref{lem:mo-compression} & \lean{DeciENSSE.Paper.lem_mo_compression} & Markers are sequences over all levels, of which only those below j are used; Exp is expandTo, validity of the image family is Admissible, and E\_j is a duplicate-free list. Inhabited supplies the default letter of the canonical hierarchy, which is unused because w and its hierarchy levels are nonempty. \\
\cref{lem:mo-closed-form} & \lean{DeciENSSE.Paper.lem_mo_closed_form} & G-sharp, L-sharp and Xi-sharp are anchoredDictionaryBound, anchoredLengthBound and anchoredExponent; N at least 1 is not needed. \\
\cref{lem:nm-infinite}(1) & \lean{DeciENSSE.Paper.lem_nm_infinite_1} & StreamPeriod also permits the trivial zero shift; the original period p remains positive. The alphabet need not be finite. \\
\cref{lem:nm-infinite}(2) & \lean{DeciENSSE.Paper.lem_nm_infinite_2} & Under the least-period hypothesis on xi, periods of suffixes are quantified directly; choosing k = |alpha| and q = |lambda| gives the lasso consequence. The suffix also retains the least period. The alphabet need not be finite. \\
\cref{lem:nm-infinite}(3) & \lean{DeciENSSE.Paper.lem_nm_infinite_3} & A default alphabet letter presents periodic streams; its value is irrelevant for nonempty periods. The alphabet need not be finite. \\
\cref{lem:nm-infinite}(4) & \lean{DeciENSSE.Paper.lem_nm_infinite_4} & Lassos are represented by their eventual periods and prefix lengths. The alphabet need not be finite. \\
\cref{lem:nm-hit}(1) & \lean{DeciENSSE.Paper.lem_nm_hit_1} & The alphabet need not be finite. \\
\cref{lem:nm-hit}(2) & \lean{DeciENSSE.Paper.lem_nm_hit_2} & The raw cut characterization is paired with the AlignedHit identity for an arbitrary relation, which includes both guard and head hit bits. The alphabet need not be finite. The automaton state carriers need not be finite. \\
\cref{lem:nm-graph-invariant} & \lean{DeciENSSE.Paper.lem_nm_graph_invariant} & PreValid packages the node identity, spine safety, target openness and the two phase conditions. Neither a finite alphabet nor (FH) is needed. E may be any finite set of states. \\
\cref{lem:nm-graph-correct} & \lean{DeciENSSE.Paper.lem_nm_graph_correct} & Neither a finite alphabet nor (FH) is needed. E may be any finite set of states. \\
\cref{lem:nm-good} & \lean{DeciENSSE.Paper.lem_nm_good} & Survivor-freeness is the conjunction of its cross and self parts. Neither a finite alphabet nor (FH) is needed. The equivalence does not require periodicity of the stream. \\
\cref{lem:nm-ends}(1) & \lean{DeciENSSE.Paper.lem_nm_ends_1} & The words are the suffixes at the two periodic cuts; their lengths are the end distances. The alphabet need not be finite. \\
\cref{lem:nm-ends}(2) & \lean{DeciENSSE.Paper.lem_nm_ends_2} & Rest distances are natural subtractions; equality also identifies periodic and rest cuts. Only the two periodic-split certificates and the suffix equality are assumed; p need not be positive or least. The alphabet need not be finite. \\
\cref{lem:nm-ends}(3) & \lean{DeciENSSE.Paper.lem_nm_ends_3} & The minimum bound is a conjunction of the two inequalities. The alphabet need not be finite. \\
\cref{lem:nm-target-end} & \lean{DeciENSSE.Paper.lem_nm_target_end} & TargetOpen packages all proper cuts; the periodic end has length t.length - rt. Neither a finite alphabet nor (FH) is needed. \\
\cref{lem:nm-rigid}(1) & \lean{DeciENSSE.Paper.lem_nm_rigid_1} & The alphabet need not be finite. \\
\cref{lem:nm-rigid}(2) & \lean{DeciENSSE.Paper.lem_nm_rigid_2} & The alphabet need not be finite. \\
\cref{lem:nm-rigid}(3) & \lean{DeciENSSE.Paper.lem_nm_rigid_3} & Uniqueness is equality of two cuts with the same non-self-similar suffix stream. The alphabet need not be finite. Cut uniqueness does not require periodicity of the appended stream. \\
\cref{lem:nm-rigid}(4) & \lean{DeciENSSE.Paper.lem_nm_rigid_4} & The alphabet need not be finite. \\
\cref{lem:nm-deletion}(1) & \lean{DeciENSSE.Paper.lem_nm_deletion_1} & The three numbered premises are jointly required in each entry. D1 uses loop equations, D2 is expanded into the five deep/shallow implications (wholly shallow pairs lift automatically), and D3 is the surviving-head disjunction. The output is the specified deleted pair. Neither a finite alphabet nor (FH) is needed. The deletion lemma holds for arbitrary streams. \\
\cref{lem:nm-deletion}(2) & \lean{DeciENSSE.Paper.lem_nm_deletion_2} & The three numbered premises are jointly required in each entry. D1 uses loop equations, D2 is expanded into the five deep/shallow implications (wholly shallow pairs lift automatically), and D3 is the surviving-head disjunction. The output is the specified deleted pair. Neither a finite alphabet nor (FH) is needed. The deletion lemma holds for arbitrary streams. \\
\cref{lem:nm-deletion}(3) & \lean{DeciENSSE.Paper.lem_nm_deletion_3} & The three numbered premises are jointly required in each entry. D1 uses loop equations, D2 is expanded into the five deep/shallow implications (wholly shallow pairs lift automatically), and D3 is the surviving-head disjunction. The output is the specified deleted pair. Neither a finite alphabet nor (FH) is needed. The deletion lemma holds for arbitrary streams. \\
\cref{lem:nm-dischargers}(1) & \lean{DeciENSSE.Paper.lem_nm_dischargers_1} & Absorption is stated for an arbitrary prefix and target stream, encompassing implication (ii). Periodicity of the stream is not needed. The alphabet need not be finite. \\
\cref{lem:nm-dischargers}(2) & \lean{DeciENSSE.Paper.lem_nm_dischargers_2} & The two impossible premises are stated directly for deep-deep and deep-shallow self pairs. The appended stream need not be periodic. The alphabet need not be finite. \\
\cref{lem:nm-four}(1) & \lean{DeciENSSE.Paper.lem_nm_four_1} & N1 is the product-state cardinality; Shrinks packages a good pair with the same endpoint and smaller total length. Neither a finite alphabet nor (FH) is needed. \\
\cref{lem:nm-four}(2) & \lean{DeciENSSE.Paper.lem_nm_four_2} & P may be any finite set of partnered rest cuts, including all of C\_u; N2 is its signature-space cardinality. Shrinks packages the conclusion. Neither a finite alphabet nor (FH) is needed. \\
\cref{lem:nm-four}(3) & \lean{DeciENSSE.Paper.lem_nm_four_3} & P may be any finite set of unpartnered spine rest cuts, including U\_u; N3 is its signature-space cardinality. Shrinks packages the conclusion. No periodic-split certificate for the target is needed. Neither a finite alphabet nor (FH) is needed. \\
\cref{lem:nm-four}(4) & \lean{DeciENSSE.Paper.lem_nm_four_4} & P may be any finite set of unpartnered target rest cuts, including U\_t; N4 is its signature-space cardinality. Shrinks packages the conclusion. No periodic-split certificate for the spine is needed. Neither a finite alphabet nor (FH) is needed. \\
\cref{lem:nm-count} & \lean{DeciENSSE.Paper.lem_nm_count} & partnered and unpartnered enumerate C and U; signature cardinalities express N1 through N4, and pumpBound is their sum plus b times the period. Neither a finite alphabet nor (FH) is needed. \\
\cref{thm:bg-nsse-sat}(1) & \lean{DeciENSSE.Paper.thm_bg_nsse_sat_1} & Variables are Fin k; finite solutions are finite terms. The cycle witness bound is k * k. \\
\cref{thm:bg-nsse-sat}(2) & \lean{DeciENSSE.Paper.thm_bg_nsse_sat_2} & Variables are Fin k, with signed variables Fin (2 * k); finite solutions are finite terms. \\
\cref{thm:bg-nsse-sat}(3) & \lean{DeciENSSE.Paper.thm_bg_nsse_sat_3} & Variables are Fin k and regular solutions are finite-graph assignments. \\
\cref{lem:bg-cores}(1) & \lean{DeciENSSE.Paper.lem_bg_cores_1} & Ordinary cores are original-word positions; symbolic cores are the iterated-option start states of a monitor. A default letter defines the total hierarchy. \\
\cref{lem:bg-cores}(2) & \lean{DeciENSSE.Paper.lem_bg_cores_2} & IsGreatest supplies the maximum ordinary core; the set identity includes the asserted comparisons. A default letter defines the total hierarchy. \\
\cref{lem:bg-exchange-types} & \lean{DeciENSSE.Paper.lem_bg_exchange_types} & ClosedType is the finite type of exchange tables; normalTypeBound is C(N,d). A default letter defines the total hierarchy. The bound applies to the entire exchange-table type for any nonempty word satisfying the depth inequalities; holehood is not needed. \\
\cref{lem:bg-comparison-map} & \lean{DeciENSSE.Paper.lem_bg_comparison_map} & NoCollapse states the predecessor-pair condition with zero-based letter indices. \\
\cref{lem:bg-level-maps} & \lean{DeciENSSE.Paper.lem_bg_level_maps} & LevelMapResult packages the hole, marker, dictionary, higher-map, start-count and unary-level claims; depth and dictionary cardinalities are also explicit. A default letter defines the total hierarchy. This structural result does not require finitely many monitor states. \\
\cref{lem:bg-isolated-exchange} & \lean{DeciENSSE.Paper.lem_bg_isolated_exchange} & The first conjunct lists the hypotheses of lem:bg-level-maps for the map identify x y: no collapse of predecessor pairs, the mapped-level hole, dictionary containment, the fixed start prefix and the unchanged run of cores. LevelMapResult includes the unchanged start counts and unary levels; letterWeight is the level-zero expansion length. A default letter defines the total hierarchy. \\
\cref{lem:bg-dictionary}(1) & \lean{DeciENSSE.Paper.lem_bg_dictionary_1} & A default letter defines the total hierarchy; Normal packages distinct exchange types of distinct isolated letters. \\
\cref{lem:bg-dictionary}(2) & \lean{DeciENSSE.Paper.lem_bg_dictionary_2} & Boundary propagation is the equivalent subtraction-free inequality b\_i + 1 ≤ b\_(i+1) + g\_(i+1); normalAllowance is A(N,d). A default letter defines the total hierarchy. \\
\cref{lem:bg-dictionary}(3) & \lean{DeciENSSE.Paper.lem_bg_dictionary_3} & normalAllowance and dictionaryBound are A and G; the final existential applies to every hole of depth at most d. A default letter defines the total hierarchy. \\
\cref{lem:bg-unary} & \lean{DeciENSSE.Paper.lem_bg_unary} &  \\
\cref{lem:bg-columns} & \lean{DeciENSSE.Paper.lem_bg_columns} & Tuples are Fin k indexed; a supported letter map expresses the replacement, whose Agree field records injectivity and preservation of all derived-monitor tables; terminal-summary bits are retained explicitly. Deletion is stated for a factorization at the two cuts. Raw tuple shortening does not require marker avoidance; the supported replacement preserves avoidance when it is assumed. \\
\cref{lem:bg-agreement} & \lean{DeciENSSE.Paper.lem_bg_agreement} & Finite supports and families are lists; the conclusion explicitly includes avoidance of the mapped marker. Agreement under derivation also holds for infinite state carriers. \\
\cref{lem:bg-compression} & \lean{DeciENSSE.Paper.lem_bg_compression} & E is a duplicate-free list of supported level-j letters; PositiveDictionaryBound, Agree and Admissible package dictionary bounds, all-state agreement with injectivity, and recursive marker validity. descentFactor N d g is col(N+d+1,g). The recursion is explicit: supports Es h with Es j = E, shortenings tau h from column shortening and level maps rho h, with rho (h+1) X = (tau h X).map (rho h) and the new marker of level h equal to the image of the old one under rho h, at every level h below j. A default letter defines the total hierarchy. \\
\cref{lem:bg-estimate} & \lean{DeciENSSE.Paper.lem_bg_estimate} & dictionaryExponent is Xi, dictionaryBound is G, and holeLengthBound N is lengthBound N d\_N; the estimates hold even at N = 0. \\
\cref{prop:ap-subsumption}(1) & \lean{DeciENSSE.Paper.prop_ap_subsumption_1} & The consequent, the conjunction of C\_1 and R\_12, is Scheme.subsumptionGoal, whose context part ranges over the common domain, which is Dom Gamma\_1 here; the class C is an arbitrary predicate on trees. For the bottom-constructor signature, instantiate C with the top-free trees of the intended class; the variable carriers need not be finite. \\
\cref{prop:ap-subsumption}(2) & \lean{DeciENSSE.Paper.prop_ap_subsumption_2} & The class C is an arbitrary predicate on trees; a solution of C\_2 assigns the free variables and the bound variables of S\_2. For the bottom-constructor signature, instantiate C with the top-free trees of the intended class; the variable carriers need not be finite. \\
\cref{cor:ap-decidable} & \lean{DeciENSSE.Paper.cor_ap_decidable} & Decidability is a computable pair of triples of the public deciders, over the full signature at every variance (arrow types are n = 2 with Examples.arrow) and over the bottom-constructor signature; the agreement of regular and arbitrary trees, for both signatures, is the propositional component. \\
\cref{ex:ap-identity} & \lean{DeciENSSE.Paper.ex_ap_identity} & Both schemes have no free variables and empty contexts; the claims hold for every predicate C on trees, without closure assumptions. The parity claim is about the flattened consequent of S\_1 ≤ S\_2. \\
\cref{ex:ap-pottier} & \lean{DeciENSSE.Paper.ex_ap_pottier} & The class D is closed under the constants and the constructor. Proving k ≤ a through the lower bound means entailing that k is below the arrow from a to the bottom tree; the signed literals are SignedC at the signed variables a+ = nu\_0 a and a- = nu\_1 a. \\
\cref{ex:ap-selfapp} & \lean{DeciENSSE.Paper.ex_ap_selfapp} & The class C is closed under the constants and the constructor; the two schemes have no free variables and empty contexts. \\
\cref{def:as-certificate}(1) & \lean{DeciENSSE.Paper.def_as_certificate_1} & Words are the type CounterexampleWord with a kind index; the entry characterizes the (I0)-words of each length L. The type does not assume the systems satisfiable. \\
\cref{def:as-certificate}(2) & \lean{DeciENSSE.Paper.def_as_certificate_2} & Words are the type CounterexampleWord with a kind index; the entry characterizes the (U)-words of each length L. The type and characterization do not assume the systems satisfiable. \\
\cref{def:as-certificate}(3) & \lean{DeciENSSE.Paper.def_as_certificate_3} & Words are the type CounterexampleWord with a kind index; the entry characterizes the hole words of each length L and includes the length |u| + |t| + 1 of the separator word u\#t. The type and characterization do not assume the systems satisfiable. \\
\cref{thm:as-main}(1) & \lean{DeciENSSE.Paper.thm_as_main_1} & Decidability is a computable triple of the public deciders, for arbitrary, regular and finite trees, on flat systems over finite ordinals; expressions are decided by decideExistsEntailsTerm and its variants (thm\_intro\_main). \\
\cref{thm:as-main}(2) & \lean{DeciENSSE.Paper.thm_as_main_2} & Decidability is a computable triple of the public deciders for the bottom-constructor signature with f covariant, over top-free arbitrary, regular and finite trees, on flat systems over finite ordinals. \\
\cref{thm:as-main}(3) & \lean{DeciENSSE.Paper.thm_as_main_3} & D ranges over the three tree classes; the observations of a word are a list of signed leaf atoms. Independence of n and c is the form of the bound. \\
\cref{thm:as-main}(4) & \lean{DeciENSSE.Paper.thm_as_main_4} & Systems over the bottom-constructor signature are flat systems read at the covariant variance with top-free values in D; the gadget pair is gadgetAnte and gadgetCons over complete lists xs and zs, with the visible carrier Sum X X and the hidden carrier Sum Z Z. \\
\cref{lem:as-countermodel}(1) & \lean{DeciENSSE.Paper.lem_as_countermodel_1} & A and O are CounterexampleWord.seeds and observations; mu\_phi and Cl\_phi are by definition the hulls of the empty set and of a singleton, and the decoding rho\_A has the atom set H\_phi(A). The satisfiability of psi is not needed. \\
\cref{lem:as-countermodel}(2) & \lean{DeciENSSE.Paper.lem_as_countermodel_2} & The satisfiability of psi is not needed; rho\_A is identified by its atom set H\_phi(A). \\
\cref{lem:as-countermodel}(3) & \lean{DeciENSSE.Paper.lem_as_countermodel_3} & B is the list CounterexampleWord.bounds of ground bounds with finite-term trees; the union of phi and B is SatT together with the bounds; no extension exists even with arbitrary values. The satisfiability of psi is not needed. \\
\cref{lem:as-countermodel}(4) & \lean{DeciENSSE.Paper.lem_as_countermodel_4} & Systems over the bottom-constructor signature are read at the covariant variance with top-free values; the gadget pair is gadgetAnte and gadgetCons over complete lists xs and zs, and the restriction is the composition of rho with Sum.inl. No counterexample word or satisfiability assumption on the gadget consequent is required; the input literals may use the full signature, interpreted on top-free values. \\
\cref{prop:as-bounds}(1) & \lean{DeciENSSE.Paper.prop_as_bounds_1} & No satisfiability hypothesis is needed. \\
\cref{prop:as-bounds}(2) & \lean{DeciENSSE.Paper.prop_as_bounds_2} & holeBound(N\_U, d\_\{N\_U\}) is holeLengthBound at the N\_U states of the unit reader, the type TargetState X Z; the satisfiability of psi is not needed. \\
\cref{prop:as-bounds}(3) & \lean{DeciENSSE.Paper.prop_as_bounds_3} & (I0) and the satisfiability of psi are not needed; Lnm and Len-sharp are extractedBoundSelf at the state counts of extractSelf, and condI0Bound is the bound of (1). This item includes the closing claims; that the mortal bound is the largest is proved for the closed forms, E3(10v+10s+17) ≤ E3(10v+12s+15), and every exact bound is at most the mortal closed form. \\
\cref{cor:as-plain} & \lean{DeciENSSE.Paper.cor_as_plain} & Plain entailment is existential entailment with no hidden variables and the consequent query x y (bfQuery over the bottom-constructor signature). Decidability is a computable pair of triples of the public deciders; the two equivalences are the propositional component, with Sat\_fin(phi) a solution by finite terms. \\
\cref{lem:bg-periodic} & \lean{DeciENSSE.Paper.lem_bg_periodic} & The infinite repetition is a periodic stream; its arbitrary default letter is immaterial for a nonempty period. \\
\cref{lem:bg-normalization}(1) & \lean{DeciENSSE.Paper.lem_bg_normalization_1} &  \\
\cref{lem:bg-normalization}(2) & \lean{DeciENSSE.Paper.lem_bg_normalization_2} & The chained equivalence is two equivalences with the same left-hand side. \\
\cref{lem:bg-normalization}(3) & \lean{DeciENSSE.Paper.lem_bg_normalization_3} &  \\
\cref{lem:bg-normalization}(4) & \lean{DeciENSSE.Paper.lem_bg_normalization_4} & Subtrees are optional, so the identity also covers words outside the domain; regularity means finitely many subtrees. \\
\cref{lem:bg-bf-order} & \lean{DeciENSSE.Paper.lem_bg_bf_order} & TopFree expresses the smaller signature; constructor positions are tested by their labels. \\
\cref{lem:bg-signed} & \lean{DeciENSSE.Paper.lem_bg_signed} & Dual coherence is the equation signDual A = A; finite and regular trees use finite-term and finite-graph presentations. Existence and uniqueness express the bijection. The class-preservation claims for decode hold without dual coherence. \\
\cref{lem:bg-flattening} & \lean{DeciENSSE.Paper.lem_bg_flattening} & C is an arbitrary tree predicate containing the values of the input subexpressions, expressed by EvalClosed; the three standard tree classes D are specializations, stated separately. Over the bottom-constructor signature at every variance, values are top-free trees of the class and the flattened systems have no top literal. Size is the number of variable occurrences (at most 6 per symbol) and of names (at most 1 per symbol). \\
\cref{lem:bg-sandwich} & \lean{DeciENSSE.Paper.lem_bg_sandwich} & Positive arity is a NeZero instance; sandwichArgs selects argument zero and repeats the filler elsewhere. \\
\cref{thm:bg-nsse}(1) & \lean{DeciENSSE.Paper.thm_bg_nsse_1} & The variable set V is finite, as in the definition of a system; finite-graph presentations express regular satisfiability. \\
\cref{thm:bg-nsse}(2) & \lean{DeciENSSE.Paper.thm_bg_nsse_2} & Decidability is a computable definition, the executable finite-satisfiability procedure Core.Mixed.decideSatFin on systems over a finite variable carrier. \\
\cref{lem:bg-extension} & \lean{DeciENSSE.Paper.lem_bg_extension} & The result holds without finiteness assumptions on the alphabet or the state carrier. \\
\cref{thm:bg-brt} & \lean{DeciENSSE.Paper.thm_bg_brt} & BoundedRejectedTail existentially selects a hole and its first rejected cut, with tail length at most the displayed bound. The alphabet need not be finite. \\
\cref{lem:bg-alignment} & \lean{DeciENSSE.Paper.lem_bg_alignment} & Strict cut inequalities are hypotheses; the two comparison predicates are expressed by their suffix-prefix conditions. \\
\cref{lem:bg-return-gap} & \lean{DeciENSSE.Paper.lem_bg_return_gap} & ReturnGap quantifies the rejected word, its first rejected cut, and the comparison with equal states. The return-gap statement needs neither finite states nor a finite alphabet. \\
\cref{thm:bg-depth} & \lean{DeciENSSE.Paper.thm_bg_depth} & J is the first rejected cut, so the tail length is w.length - J; a default letter defines the total canonical hierarchy. The alphabet need not be finite. \\
\cref{thm:bg-search} & \lean{DeciENSSE.Paper.thm_bg_search} & InL d expresses depth at most d; lettered holes are nonempty by definition. Both monitor and reader forms are included; a default letter defines the total hierarchy. The reader version, like the monitor version, does not require a finite alphabet. \\
\cref{cor:bg-hole-decision} & \lean{DeciENSSE.Paper.cor_bg_hole_decision} & holeLengthBound N is lengthBound N (2 * N * (4 * N - 1) - 1). \\
\cref{lem:bg-desub} & \lean{DeciENSSE.Paper.lem_bg_desub} & The one-step equivalence is oriented from derivative to expansion; the canonical hierarchy uses a default letter. Cores Q j is the iterated option state type. \\
\cref{thm:ch-characterization}(1) & \lean{DeciENSSE.Paper.thm_ch_characterization_1} & Visible and hidden carriers need not be finite. \\
\cref{thm:ch-characterization}(2) & \lean{DeciENSSE.Paper.thm_ch_characterization_2} & Includes the unnumbered satisfiability-dispatch equivalence. Visible and hidden carriers need not be finite. \\
\cref{lem:ch-pure}(1) & \lean{DeciENSSE.Paper.lem_ch_pure_1} & Visible and hidden carriers need not be finite. \\
\cref{lem:ch-pure}(2) & \lean{DeciENSSE.Paper.lem_ch_pure_2} & Visible and hidden carriers need not be finite. \\
\cref{lem:tc-extension} & \lean{DeciENSSE.Paper.lem_tc_extension} & Consistency separates forbidden roots from exclusion pairs; visible equality is inclusion because seeds are retained. The decoded extension is also least in the information order. Visible and hidden carriers may be infinite. \\
\cref{lem:tc-two-seeds} & \lean{DeciENSSE.Paper.lem_tc_two_seeds} & The asserted and denied sets are a list of at most two signed observations. The result also holds for infinite visible and hidden carriers. \\
\cref{lem:tc-pins} & \lean{DeciENSSE.Paper.lem_tc_pins} & The four covariant cases are packaged by Observation.bound; signed observations package the mixed cases. Bounds contain finite terms. \\
\cref{prop:ch-meet}(1) & \lean{DeciENSSE.Paper.prop_ch_meet_1} & Models and projected models are sets of signed atoms; their intersection closure packages the assignment-level consequence. Intersection closure also holds over arbitrary variable carriers. \\
\cref{prop:ch-meet}(2) & \lean{DeciENSSE.Paper.prop_ch_meet_2} & Positive arity is witnessed by a child index; infinitude is absence of a finite-term presentation. \\
\cref{lem:tc-pinning-regular} & \lean{DeciENSSE.Paper.lem_tc_pinning_regular} & The variable set V is finite, as in the definition of a system; regular bounds are bindings to a finite indexed family of finite graphs. \\
\cref{thm:tc-pinning} & \lean{DeciENSSE.Paper.thm_tc_pinning} & Bounds are a finite list of finite-term comparisons; finiteness means having a finite-term presentation. \\
\cref{lem:tc-signed}(1) & \lean{DeciENSSE.Paper.lem_tc_signed_1} & System satisfaction and satisfaction of the bound list are paired as a conjunction. \\
\cref{lem:tc-signed}(2) & \lean{DeciENSSE.Paper.lem_tc_signed_2} &  \\
\cref{lem:tc-signed}(3) & \lean{DeciENSSE.Paper.lem_tc_signed_3} & Finiteness is existence of finite-term presentations; the inverse is the library function decode. Finiteness of decode does not require dual coherence. \\
\cref{lem:tc-median}(1) & \lean{DeciENSSE.Paper.lem_tc_median_1} & Monotonicity in the three arguments is packaged as simultaneous monotonicity. \\
\cref{lem:tc-median}(2) & \lean{DeciENSSE.Paper.lem_tc_median_2} &  \\
\cref{lem:tc-median}(3) & \lean{DeciENSSE.Paper.lem_tc_median_3} & Finite trees are trees with a finite domain. \\
\cref{lem:tc-median}(4) & \lean{DeciENSSE.Paper.lem_tc_median_4} & System closure and closure of arbitrary lower and upper ground bounds are packaged separately. \\
\cref{lem:tc-pinning-covariant} & \lean{DeciENSSE.Paper.lem_tc_pinning_covariant} & Bounds are a finite list of finite-term comparisons. \\
\cref{thm:tc-regular}(1) & \lean{DeciENSSE.Paper.thm_tc_regular_1} & The variable set V is finite, as in the definition of a system; regularity is existence of a finite-graph presentation. \\
\cref{thm:tc-regular}(2) & \lean{DeciENSSE.Paper.thm_tc_regular_2} & ClassExtension quantifies the regular visible assignment and its extension. \\
\cref{thm:tc-regular}(3) & \lean{DeciENSSE.Paper.thm_tc_regular_3} & This item also packages the theorem's unnumbered entailment equivalence; regularity uses finite-graph presentations. \\
\cref{thm:tc-finite}(1) & \lean{DeciENSSE.Paper.thm_tc_finite_1} & Finiteness is existence of a finite-term presentation. Visible and hidden carriers may be infinite. \\
\cref{thm:tc-finite}(2) & \lean{DeciENSSE.Paper.thm_tc_finite_2} & ClassExtension packages extension of finite assignments; this item also includes the theorem's unnumbered finite dispatch equivalence. Only the visible carrier is required to be finite. \\
\cref{cor:tc-decidable} & \lean{DeciENSSE.Paper.cor_tc_decidable} & The conditional is a computable transformer at fixed finite carriers. A decision procedure over arbitrary trees yields one over regular trees through the regular equivalence, and one over finite trees through the finite dispatch, whose two finite-satisfiability tests run Core.Mixed.decideSatFin. Both equivalences are included. \\
\cref{lem:tc-gadget}(1) & \lean{DeciENSSE.Paper.lem_tc_gadget_1} &  \\
\cref{lem:tc-gadget}(2) & \lean{DeciENSSE.Paper.lem_tc_gadget_2} & Omega is the unfolding of the one-state graph. Finite trees use finite terms; height counts constructors (at arity zero the witness is still f()). \\
\cref{thm:tc-gadget} & \lean{DeciENSSE.Paper.thm_tc_gadget} & The three tree classes are a conjunction. Fresh variables are sum copies; lists enumerate all visible and hidden variables. Top-free values interpret the smaller signature. The equivalences also allow full-syntax systems interpreted with top-free values; the paper uses their specialization without top literals. \\
\cref{cor:tc-bf} & \lean{DeciENSSE.Paper.cor_tc_bf} & Fin-indexed public syntax, with k visible and h hidden variables. Each conditional is a computable transformer: a decision procedure for the covariant problem over the full signature, on the doubled carriers, yields one for the bottom-constructor signature in the same tree class, through the reduction toCoreTopFree and splitTopFree. The three reductions and the regular/arbitrary equivalence are included. The size claims are the doubled carriers Sum (Fin k) (Fin k) and Sum (Fin h) (Fin h) and the translated lengths, two more literals per visible or hidden variable. \\
\cref{lem:ex-head-before-J} & \lean{DeciENSSE.Paper.lem_ex_head_before_J} & J is an option; its defined value is supplied as an equality. The alphabet need not be finite. \\
\cref{lem:ex-fh} & \lean{DeciENSSE.Paper.lem_ex_fh} &  \\
\cref{thm:ex-reduction}(1) & \lean{DeciENSSE.Paper.thm_ex_reduction_1} & This item includes the theorem's unnumbered pointwise equivalence; signed variables are pairs and visible variables use Sum.inl. \\
\cref{thm:ex-reduction}(2) & \lean{DeciENSSE.Paper.thm_ex_reduction_2} & Models are sets of signed atoms; `restrict` restricts a model to the visible variables and `configSet` is the atom set of an assignment. \\
\cref{cor:ex-entailment} & \lean{DeciENSSE.Paper.cor_ex_entailment} &  \\
\cref{ex:intro-shift} & \lean{DeciENSSE.Paper.ex_intro_shift} & The running systems use top-free binary trees at the covariant order, with an unused visible b in phi\_1 and phi\_3. The subtree-position identity allows arbitrary trees; the subtree characterization requires only a top-free x; the explicit witness theorem for phi\_2 also holds over the full signature. Kpos and Suf are sets of words, the shifted subtree is subtreeBot, and the nested bound of phi\_3 is stated semantically. \\
\cref{thm:intro-main} & \lean{DeciENSSE.Paper.thm_intro_main} & Decidability is a computable pair of triples of the public deciders for systems of type expressions, over the full signature at every variance and over the bottom-constructor signature with f covariant. The certificate claims, for flat systems in each tree class D, are the propositional component; E3(O(s)) is E3(22s + 15), and E3(44s + 15) over the bottom-constructor signature, where the word is one of the gadget pair gadgetAnte and gadgetCons over complete lists xs and zs. \\
\cref{thm:lf-representation}(1) & \lean{DeciENSSE.Paper.thm_lf_representation_1} & LeafPair packages disjointness and closure under arbitrary right extensions. \\
\cref{thm:lf-representation}(2) & \lean{DeciENSSE.Paper.thm_lf_representation_2} &  \\
\cref{thm:lf-representation}(3) & \lean{DeciENSSE.Paper.thm_lf_representation_3} & The four constant-cone equalities are expressed by membership for either sign. \\
\cref{lem:lf-clauses} & \lean{DeciENSSE.Paper.lem_lf_clauses} & The singleton-system clause equivalence includes the automatic extension and exclusion conditions; uniqueness and existence express the bijection. The correspondence also holds for infinite variable carriers. \\
\cref{lem:lf-moves}(1) & \lean{DeciENSSE.Paper.lem_lf_moves_1} & Includes the stated extension to pure derivations. The variable carrier need not be finite. \\
\cref{lem:lf-moves}(2) & \lean{DeciENSSE.Paper.lem_lf_moves_2} & Includes the stated sign invariant for pure derivations. The variable carrier need not be finite. \\
\cref{lem:lf-closure}(1) & \lean{DeciENSSE.Paper.lem_lf_closure_1} & Facts, seeds, closure and leastness are separate conjuncts. The variable carrier need not be finite. \\
\cref{lem:lf-closure}(2) & \lean{DeciENSSE.Paper.lem_lf_closure_2} & Consistency is Avoids paired with Consistent; the general leastness clause also applies to the empty and singleton seed sets. The variable carrier need not be finite. \\
\cref{lem:lf-closure}(3) & \lean{DeciENSSE.Paper.lem_lf_closure_3} & Consistency separates forbidden roots from exclusion pairs; Exclusion fixes the positive-first orientation. The variable carrier need not be finite. \\
\cref{lem:lf-closure}(4) & \lean{DeciENSSE.Paper.lem_lf_closure_4} & The variable carrier need not be finite. \\
\cref{lem:lf-closure}(5) & \lean{DeciENSSE.Paper.lem_lf_closure_5} & Sum types separate visible and hidden variables; consistency is Avoids paired with Consistent. The equivalence holds for every seed set, without the paper's redundant closure assumptions. Visible and hidden carriers may be infinite. \\
\cref{lem:lf-duality} & \lean{DeciENSSE.Paper.lem_lf_duality} & The variable carrier need not be finite. \\
\cref{lem:lf-ext-first} & \lean{DeciENSSE.Paper.lem_lf_ext_first} & The variable carrier need not be finite. \\
\cref{lem:lf-extensions} & \lean{DeciENSSE.Paper.lem_lf_extensions} & The variable carrier need not be finite. \\
\cref{lem:lf-collision} & \lean{DeciENSSE.Paper.lem_lf_collision} & The variable carrier need not be finite. \\
\cref{cor:lf-dead}(1) & \lean{DeciENSSE.Paper.cor_lf_dead_1} & This item states sufficiency of cause (a) and its pure form; the full disjunction equivalence is in item (3). The variable carrier need not be finite. \\
\cref{cor:lf-dead}(2) & \lean{DeciENSSE.Paper.cor_lf_dead_2} & This item states sufficiency of cause (b); the full disjunction equivalence is in item (3). The variable carrier need not be finite. \\
\cref{cor:lf-dead}(3) & \lean{DeciENSSE.Paper.cor_lf_dead_3} & Includes the full three-cause equivalence and the pure form of cause (c). The variable carrier need not be finite. \\
\cref{lem:mo-sep}(1) & \lean{DeciENSSE.Paper.lem_mo_sep_1} & The separator word is sepWord u t over Option alpha, with none as the separator. Alphabet finiteness is not required. \\
\cref{lem:mo-sep}(2) & \lean{DeciENSSE.Paper.lem_mo_sep_2} & Cross comparisons are TwoWordInstance.Cross, which also records the cut bounds. Alphabet finiteness is not required. \\
\cref{lem:mo-sep}(3) & \lean{DeciENSSE.Paper.lem_mo_sep_3} & The comparison equivalence holds over any alphabet; finiteness is not required. \\
\cref{lem:mo-sep}(4) & \lean{DeciENSSE.Paper.lem_mo_sep_4} & An anchored comparison is a comparison whose continuation letter, at the cut x + length w - y, is the separator; aligned cross comparisons are TwoWordInstance.Aligned. Alphabet finiteness is not required. \\
\cref{lem:mo-pi-reader}(1) & \lean{DeciENSSE.Paper.lem_mo_pi_reader_1} & The reader is Mortal.reader on all states of RState, itself a computable definition. This item also packages the computability of the relations as a computable definition, whose values are the decision procedures Self.instDecidableAdm, Mortal.instDecidableRej and Mortal.instDecidablePi for admission, rejection and Pi. \\
\cref{lem:mo-pi-reader}(2) & \lean{DeciENSSE.Paper.lem_mo_pi_reader_2} & A Pi-hole is written out as a nonempty reader hole with an anchored comparison in Pi; the admissions are Self.Adm, the cross and self guards together. \\
\cref{lem:mo-pi-reader}(3) & \lean{DeciENSSE.Paper.lem_mo_pi_reader_3} & Stated for the mortal hole (u, t) of item (2); the rejected tail is the number of letters after the first rejected cut. \\
\cref{lem:mo-depth} & \lean{DeciENSSE.Paper.lem_mo_depth} & The depth bound d\_I is written out with the state count of the reader. \\
\cref{lem:mo-base} & \lean{DeciENSSE.Paper.lem_mo_base} & The lettered monitor of a reader is ofReader and the base channel is baseChan, whose relation is also false into the start label, which no comparison reaches; the first claim holds for every word; a Pi-hole is written out as a nonempty reader hole with an anchored comparison in Pi. Neither alphabet nor state finiteness is required. \\
\cref{lem:mo-desub} & \lean{DeciENSSE.Paper.lem_mo_desub} & The expansion of v' is expand l v'; the derived monitor and channel are der and posDer. State finiteness is not required. \\
\cref{cor:mo-tower} & \lean{DeciENSSE.Paper.cor_mo_tower} & The canonical hierarchy is hierOf with markers markOf, which needs an inhabited alphabet; the towers are tower and posTower, and their recursion is included; the number of states is not used. \\
\cref{thm:mo-anchored} & \lean{DeciENSSE.Paper.thm_mo_anchored} & L-sharp, G-sharp and Xi-sharp are anchoredLengthBound, anchoredDictionaryBound and anchoredExponent. Alphabets are inhabited for the canonical hierarchy; finiteness of the alphabet and N at least 1 are not needed. A Pi-hole is written out as in lem:mo-base. \\
\cref{thm:mo-short} & \lean{DeciENSSE.Paper.thm_mo_short} & L-sharp of I is anchoredLengthBound at the state count of the reader and the depth bound of lem:mo-depth, written out. \\
\cref{cor:mo-decidable} & \lean{DeciENSSE.Paper.cor_mo_decidable} & Decidability is a computable definition, the executable procedure Self.decideHasMortalHole, together with the bounded search over words of length at most L-sharp of I that it performs. \\
\cref{thm:mo-decide} & \lean{DeciENSSE.Paper.thm_mo_decide} & Lnm is pumpHoleBound and L-sharp of I is written out as in thm:mo-short. Decidability is a computable definition for instances with (FH): the executable bounded search SelfTwoWordInstance.decideBoundedHole at the bound selfHoleBound, transported along hasHole\_iff\_bounded. \\
\cref{lem:mo-estimate} & \lean{DeciENSSE.Paper.lem_mo_estimate} & N-bar is the cardinality of the reader-state type Mortal.State, and Len-sharp is mortalHoleBound at the four cardinalities; this numerical estimate requires no automaton instance and also permits empty state carriers. \\
\cref{cor:mo-extracted} & \lean{DeciENSSE.Paper.cor_mo_extracted} & Len-sharp of I\_pq is extractedBoundSelf mortalHoleBound X Z, the mortal bound at the state counts of extractSelf, the same for all p and q; the separator word u\#t has length |u| + |t| + 1. \\
\cref{lem:nm-continuation}(1) & \lean{DeciENSSE.Paper.lem_nm_continuation_1} & State preservation is quantified over an arbitrary automaton, and thus includes all four readers on either word. Neither a finite alphabet nor (FH) is needed. \\
\cref{lem:nm-continuation}(2) & \lean{DeciENSSE.Paper.lem_nm_continuation_2} & Cross includes both cut bounds, including s ≤ |u|; self comparisons specialize the same identity to u = t and strict cuts. Neither a finite alphabet nor (FH) is needed. \\
\cref{lem:nm-continuation}(3) & \lean{DeciENSSE.Paper.lem_nm_continuation_3} & Keeps is defined by the overhang prefix test; self comparisons specialize cross comparisons to equal words. The periodic stream of a nonempty overhang takes its default letter from the overhang itself, so the alphabet may be empty; it need not be finite. \\
\cref{lem:nm-decomposition} & \lean{DeciENSSE.Paper.lem_nm_decomposition} & Completion is split into cross and self parts; the rejected cut is the length of the pre-hole target, whose J equals that length by definition. The alphabet need not be finite. \\
\cref{lem:nm-breaker} & \lean{DeciENSSE.Paper.lem_nm_breaker} & The periodic prefix test is written θ ≤prefix oθ, equivalent to θ ≤prefix o\textasciicircum{}ω for nonempty o; completion has cross and self parts. Neither a finite alphabet nor (FH) is needed. \\
\cref{lem:nm-thin}(1) & \lean{DeciENSSE.Paper.lem_nm_thin_1} & Neither a finite alphabet nor (FH) is needed. \\
\cref{lem:nm-thin}(2) & \lean{DeciENSSE.Paper.lem_nm_thin_2} & Neither a finite alphabet nor (FH) is needed. \\
\cref{lem:nm-thin}(3) & \lean{DeciENSSE.Paper.lem_nm_thin_3} & A default alphabet letter presents periodic streams; its value is irrelevant for nonempty periods. Neither a finite alphabet nor (FH) is needed. \\
\cref{lem:nm-thin}(4) & \lean{DeciENSSE.Paper.lem_nm_thin_4} & Neither a finite alphabet nor (FH) is needed. \\
\cref{lem:nm-survival}(1) & \lean{DeciENSSE.Paper.lem_nm_survival_1} & Non-alignment is equivalent to a nonempty overhang for a comparison; self comparisons are the strict equal-word specialization. A default alphabet letter presents periodic streams; its value is irrelevant for nonempty periods. The alphabet need not be finite. \\
\cref{lem:nm-survival}(2) & \lean{DeciENSSE.Paper.lem_nm_survival_2} & Self comparisons are the specialization u = t and s < e. The alphabet need not be finite. \\
\cref{lem:nm-survival}(3) & \lean{DeciENSSE.Paper.lem_nm_survival_3} & The overhang length is offset; self comparisons are the specialization u = t and s < e. A default alphabet letter presents periodic streams; its value is irrelevant for nonempty periods. The alphabet need not be finite. \\
\cref{cor:nm-lasso-completion} & \lean{DeciENSSE.Paper.cor_nm_lasso_completion} & Completion is split into cross completion and self completion. A default alphabet letter presents periodic streams; its value is irrelevant for nonempty periods. Neither a finite alphabet nor (FH) is needed. The completion implication does not require a pre-hole hypothesis. \\
\cref{prop:nm-short-prehole} & \lean{DeciENSSE.Paper.prop_nm_short_prehole} & E is a finite set of states; the result holds without requiring its members to be safe. Decidability is a computable definition, the executable bounded search TwoWordInstance.preHoleExistsDecidable over pairs of words of length at most 12aba'b'. (FH) is not needed. \\
\cref{lem:nm-aperiodic} & \lean{DeciENSSE.Paper.lem_nm_aperiodic} & Self comparisons specialize cross comparisons to equal words; survivor-freeness has cross and self parts. Neither a finite alphabet nor (FH) is needed. \\
\cref{thm:nm-pump} & \lean{DeciENSSE.Paper.thm_nm_pump} & pumpBound is the displayed polynomial and pumpCapacity is Pcap; survivor-freeness has cross and self parts. Neither a finite alphabet nor (FH) is needed. \\
\cref{thm:nm-three-cases}(1) & \lean{DeciENSSE.Paper.thm_nm_three_cases_1} & State carriers are finite initial segments; pumpHoleBound is Lnm. This item gives branching completion with its bound, without (FH). \\
\cref{thm:nm-three-cases}(2) & \lean{DeciENSSE.Paper.thm_nm_three_cases_2} & State carriers are finite initial segments; survivor-freeness has cross and self parts. Includes the exact prefix claim for a completed hole and the Lnm bound without (FH). \\
\cref{thm:nm-three-cases}(3) & \lean{DeciENSSE.Paper.thm_nm_three_cases_3} & NonMortalCase is the disjunction of the branching and thin survivor-free pre-hole cases. Includes the theorem's full equivalence and its original-hole mortal alternative. The alphabet need not be finite. \\
\cref{cor:nm-short-holes} & \lean{DeciENSSE.Paper.cor_nm_short_holes} & State carriers are finite initial segments and pumpHoleBound is Lnm. \\
\cref{lem:sm-summaries}(1) & \lean{DeciENSSE.Paper.lem_sm_summaries_1} & Pure moves are represented in the covariant signed system. \\
\cref{lem:sm-summaries}(2) & \lean{DeciENSSE.Paper.lem_sm_summaries_2} & Pure moves are represented in the covariant signed system. \\
\cref{lem:sm-summaries}(3) & \lean{DeciENSSE.Paper.lem_sm_summaries_3} & Pure moves are represented in the covariant signed system. \\
\cref{lem:sm-lowest-point} & \lean{DeciENSSE.Paper.lem_sm_lowest_point} &  \\
\cref{lem:sm-duality}(1) & \lean{DeciENSSE.Paper.lem_sm_duality_1} &  \\
\cref{lem:sm-duality}(2) & \lean{DeciENSSE.Paper.lem_sm_duality_2} &  \\
\cref{lem:sm-duality}(3) & \lean{DeciENSSE.Paper.lem_sm_duality_3} &  \\
\cref{lem:sm-mu} & \lean{DeciENSSE.Paper.lem_sm_mu} & The bound 4\textasciicircum{}|V| is the cardinality of the signed powerset; PF is the finite set defined by this bounded search, and it equals the set of the paper, defined with arbitrary words. The computability claim is a computable definition: membership in that set is decided by membership in PF. \\
\cref{lem:sm-roots}(1) & \lean{DeciENSSE.Paper.lem_sm_roots_1} & Root deadness is expanded as reachability of a forbidden root. \\
\cref{lem:sm-roots}(2) & \lean{DeciENSSE.Paper.lem_sm_roots_2} & Root deadness is expanded as reachability of a forbidden root. \\
\cref{lem:sm-roots}(3) & \lean{DeciENSSE.Paper.lem_sm_roots_3} & Root deadness is expanded as reachability of a forbidden root. \\
\cref{lem:sm-roots}(4) & \lean{DeciENSSE.Paper.lem_sm_roots_4} & Root deadness is expanded as reachability of a forbidden root. \\
\cref{lem:sm-dead-states}(1) & \lean{DeciENSSE.Paper.lem_sm_dead_states_1} & The root-dead region is represented by the computable state predicate negReachS rather than an intersection with a set of variables. \\
\cref{lem:sm-dead-states}(2) & \lean{DeciENSSE.Paper.lem_sm_dead_states_2} & flipHitS is intersection of the flipped summary with PF. \\
\cref{lem:sm-dead-states}(3) & \lean{DeciENSSE.Paper.lem_sm_dead_states_3} & The support-dead region is the disjunction of negReachS and muReachS on the current state; the last disjunct is opposite-atom membership in Mu. \\
\cref{lem:sm-dead-states}(4) & \lean{DeciENSSE.Paper.lem_sm_dead_states_4} &  \\
\cref{lem:sm-dead-states}(5) & \lean{DeciENSSE.Paper.lem_sm_dead_states_5} & The computability claim for the two finite sets of signed variables is a computable definition: membership is decided by the executable procedures decideRootDeadEmpty and decideSupportDeadEmpty; root deadness is reachability of a forbidden root. \\
\cref{thm:sm-dead} & \lean{DeciENSSE.Paper.thm_sm_dead} &  \\
\cref{lem:sm-membership} & \lean{DeciENSSE.Paper.lem_sm_membership} &  \\
\cref{lem:sm-head-prefix} & \lean{DeciENSSE.Paper.lem_sm_head_prefix} & Root deadness is expanded as reachability of a forbidden root. \\
\cref{prop:sm-decide}(1) & \lean{DeciENSSE.Paper.prop_sm_decide_1} & Forbidden roots are enumerated by negList; MuCollision is simultaneous membership of both signs. The bound is the cardinality of the product of the two flagged powerset readers, equal to 2\textasciicircum{}(4|V|+2). Decidability of satisfiability is a computable definition, the executable procedure decideSatMixed. \\
\cref{prop:sm-decide}(2) & \lean{DeciENSSE.Paper.prop_sm_decide_2} & FailureI0 is consequent membership and antecedent nonmembership at a visible signed atom. The bound is the cardinality of the product of the two flagged powerset readers, equal to 2\textasciicircum{}(2|X|+2|X+Z|+2). Decidability of (I0) is a computable definition, the executable procedure decideCondI0. \\
\cref{lem:un-dead} & \lean{DeciENSSE.Paper.lem_un_dead} & Absence of a live root-dead seed is written as a universally quantified implication. \\
\cref{lem:un-source}(1) & \lean{DeciENSSE.Paper.lem_un_source_1} & The Boolean component is specified by its truth predicate; the state count is the cardinality of the product carrier. \\
\cref{lem:un-source}(2) & \lean{DeciENSSE.Paper.lem_un_source_2} &  \\
\cref{lem:un-source}(3) & \lean{DeciENSSE.Paper.lem_un_source_3} & The comparison includes its strict cut bounds explicitly. \\
\cref{lem:un-holes} & \lean{DeciENSSE.Paper.lem_un_holes} & Includes the consequent universal characterization of (U); signed variables index the readers. \\
\cref{lem:un-rejected-path} & \lean{DeciENSSE.Paper.lem_un_rejected_path} & RejectedPath packages betweenness and the admission cone. \\
\cref{thm:un-decidable} & \lean{DeciENSSE.Paper.thm_un_decidable} & N\_U is the cardinality of TargetState; holeLengthBound N is holeBound(N, d\_N). Decidability is a computable definition for satisfiable antecedents, the executable procedure decideCondU. The closed form is stated for a nonempty visible carrier, with v = |X| and s = |X| + |Z|. \\
